\subsection{Interpretation of the Main Findings}

The experimental results reveal a clear separation between anomaly screening and conditional fault-domain diagnosis. Once an anomalous window is provided, the four annotated fault domains remain highly distinguishable even under complete flight-level separation. In contrast, anomaly detection is substantially more sensitive to the evaluation protocol. This difference indicates that the main challenge in transferring the complete diagnostic pipeline to unseen flights lies in determining whether a telemetry pattern is abnormal, rather than in distinguishing among the labelled fault domains after an anomaly has been identified.

This behavior can be explained by the different learning conditions of the two stages. Stage 1 must separate normal and anomalous operation across heterogeneous flights, where vehicle state, trajectory, operating regime, and local telemetry statistics may vary considerably between logs. Chronological splitting can retain part of this flight-specific context across partitions, whereas leave-log-out evaluation removes this shared context entirely. Stage 2 operates on a narrower conditional problem because only explicitly labelled anomalous windows are considered. The comparatively stable fault-domain performance therefore suggests that domain-specific telemetry patterns remain learnable across flights, while the boundary between normal and abnormal operation is more sensitive to flight-to-flight variation.

The comparison with MTCL-UAV further supports this interpretation from a different modeling perspective. Under the same leave-log-out data isolation, the normal-only reconstruction model achieved lower anomaly-detection performance than the supervised Stage 1 detector. MTCL-UAV explicitly models multiscale temporal dynamics, whereas HS-AeroTS relies on temporal-statistical representations followed by supervised discrimination. The result therefore suggests that increasing temporal-model complexity alone does not necessarily resolve cross-flight anomaly separation on UAV-SEAD. At the same time, the two methods use different supervision paradigms, so the comparison should be interpreted as evidence about their practical behavior under a common evaluation protocol rather than as proof of architectural superiority.

The complete five-class results provide a complementary perspective on the hierarchical design. The Cascade inherits both false negatives and false positives from Stage 1 before applying fault-domain diagnosis, and its relative performance with respect to Direct Five-Class classification changes across evaluation protocols. The hierarchical formulation should therefore not be viewed primarily as a mechanism for maximizing five-class predictive accuracy. Its main value is structural: it separates anomaly screening from fault interpretation, makes detector errors explicit, and enables class-specific diagnostic evidence to be generated only after the anomaly gate is activated. This organization is particularly important for the subsequent evidence and software-inspection stages, where the provenance of each diagnostic decision must remain traceable.

Taken together, these findings suggest that the current framework is strongest at conditional diagnosis and evidence organization, while cross-flight anomaly screening remains the principal bottleneck in the end-to-end chain. Improving the robustness of Stage 1 to flight-level distribution shifts is therefore likely to provide greater benefit to the complete system than further increasing the complexity of the downstream fault-domain classifier.
